(1) \(dS=d(Mc^2)/T_H=(8\pi Gk_B/(\hbar c))M\,dM\)，积分得
\[
S(M)=\frac{4\pi k_BG M^2}{\hbar c}
=\frac{k_Bc^3A_H}{4G\hbar}.
\]
熵与视界面积成正比而非体积；积分常数由给定参考条件固定。
(2) \(dT_H/dM=-\hbar c^3/(8\pi Gk_BM^2)\)，因此
\[
C=\frac{d(Mc^2)}{dT_H}=-\frac{8\pi Gk_BM^2}{\hbar c}<0.
\]
黑洞失去质量时 \(T_H\propto M^{-1}\) 上升。孤立自引力体系的能量与熵约束不同于普通短程相互作用正则热浴；负热容说明与无限恒温浴接触不稳定，不等于违反第一定律。
(3) 代入面积、温度和给定Stefan常数：
\[
P=\epsilon\frac{\hbar c^6}{15360\pi G^2M^2},\qquad
\frac{dM}{dt}=-\frac{P}{c^2}=-\frac{\epsilon\hbar c^4}{15360\pi G^2M^2}=-\frac{D}{M^2}.
\]
积分 \(M(t)^3=M_0^3-3Dt\)，故 \(t(M)=\frac{M_0^3-M^3}{3D}\)。
(4) 形式蒸发寿命
\[
t_{\rm evap}=\frac{M_0^3}{3D}=\frac{5120\pi G^2M_0^3}{\epsilon\hbar c^4},
\quad T_H(t)=\frac{\hbar c^3}{8\pi Gk_B(M_0^3-3Dt)^{1/3}}.
\]
临近形式终点 \(T_H\propto(t_{\rm evap}-t)^{-1/3}\)，\([G^2M^3/(\hbar c^4)]=t\)。实际辐射存在频率依赖灰体因子与粒子物种变化，且 \(M\) 逼近Planck质量时半经典近似失效，不能断言温度真实发散。
